\documentclass[10pt,twocolumn]{article}
\usepackage[margin=0.75in]{geometry}
\usepackage{times}
\usepackage{amsmath,amssymb}
\usepackage{graphicx}
\usepackage{hyperref}
\usepackage{titlesec}
\usepackage{abstract}
\usepackage{enumitem}
\usepackage{fancyhdr}

\titleformat{\section}{\normalfont\large\bfseries}{\thesection}{1em}{}
\titleformat{\subsection}{\normalfont\bfseries}{\thesubsection}{1em}{}
\setlength{\columnsep}{0.3in}

\pagestyle{fancy}
\fancyhf{}
\renewcommand{\headrulewidth}{0pt}
\fancyfoot[C]{\thepage}

\title{\vspace{-0.5em}\LARGE\bfseries REPLACE: Paper Title\\
\large (e.g. ``Systems Friction as Signal: Temporary High-Agency
Outsiders in Civic Digital Design'')}

\author{
Jade Zhao\\
\small Luddy School of Informatics, Computing, and Engineering\\
\small Indiana University Bloomington\\
\small \texttt{jlzhao@iu.edu}
}

\date{}

\begin{document}
\twocolumn[
  \begin{@twocolumnfalse}
    \maketitle
    \begin{abstract}
    \noindent
    REPLACE: 150--250 word abstract. State the problem (what friction /
    accessibility / ethics gap this addresses), the setting (Madrid study
    abroad friction, or the local-digital nonprofit context), the approach
    taken, and the finding or artifact produced. Write it last, after the
    rest of the paper is drafted -- it should be a compressed version of
    the whole paper, not a teaser.
    \end{abstract}
    \vspace{1em}
  \end{@twocolumnfalse}
]

\section{Introduction}
REPLACE: Motivate the problem in 3--4 paragraphs.
\begin{itemize}[leftmargin=*]
  \item Paragraph 1: the concrete situation that surfaced the problem
        (e.g., a specific moment of friction encountered abroad, or a
        specific nonprofit partner's constraint).
  \item Paragraph 2: why existing approaches / systems don't handle this
        well (the gap).
  \item Paragraph 3: this paper's contribution, stated as 2--3 explicit
        bullet-style claims (can be inline prose, but should be
        scannable).
  \item Paragraph 4 (optional): roadmap of the rest of the paper.
\end{itemize}

\section{Related Work}
REPLACE: Group prior work into 2--3 thematic clusters rather than a flat
list. For this material, likely clusters: (1) accessibility / WCAG
(Web Content Accessibility Guidelines) research, (2) AI ethics and
fairness frameworks, (3) civic tech / public-interest technology
literature. Cite each source in your own words -- do not quote
directly. One or two sentences per cluster is enough at draft stage.

\section{Approach / System}
REPLACE: This is the technical or methodological core. For
local-digital or madrid-ai-ethics, this section should describe:
\begin{itemize}[leftmargin=*]
  \item What was built, observed, or analyzed
  \item The design principles or ethical framework applied
  \item Concrete artifacts (screenshots, decision frameworks, code
        structure) -- figures go here
\end{itemize}

\subsection{Design Principles}
REPLACE: e.g., plain-language defaults, recoverable errors, no assumed
prior belonging -- whatever principles actually governed the work.

\section{Findings / Evaluation}
REPLACE: What happened when this was applied. For a reflective /
qualitative piece (most likely fit for Madrid material), this can be
structured as observed patterns rather than quantitative results --
still present them with the same rigor: what was observed, how it was
recorded, what it implies.

\section{Discussion}
REPLACE: Limitations, scope of applicability, what a reader should not
conclude from this work. One paragraph minimum -- reviewers (and
hiring readers) weight self-aware limitations heavily.

\section{Conclusion}
REPLACE: 3--5 sentences. Restate the contribution, not a new claim.

\section*{Acknowledgments}
REPLACE or delete. If this draws on ServeIT / Serve-AI / PIT-UN
(Public Interest Technology University Network) materials, credit them
here per the CC BY-SA 4.0 terms already governing that work.

\begin{thebibliography}{9}
\bibitem{ref1} REPLACE: Author(s), \textit{Title}, Venue, Year.
\bibitem{ref2} REPLACE: Author(s), \textit{Title}, Venue, Year.
\end{thebibliography}

\end{document}
